% =====================================================================
%  Report 6
%  A user interface for the curation of EarthPortal
%  Build: latexmk -pdf emmanuel_dataTerra_Report6-Curation-Interface.tex
%  Numbers: snapshot of 16 October 2026 (SA_Management v0.3.0, notebook 05).
%  Screenshots: production submission form of EarthPortal, 29 October 2026.
% =====================================================================
\documentclass[11pt,a4paper]{article}
\usepackage{dataterra-report}
\usetikzlibrary{positioning}

\reportsubtitle{A user interface for the curation of EarthPortal}

\begin{document}
\makereporttitlepage

% =====================================================================
\section{Introduction}

In this short report, we suggest a user interface that could facilitate the curation job on
EarthPortal. It builds on the audits produced by our pipeline (SA\_Management, v0.3.0), which give, for
each ontology, its FAIRness according to O'FAIRe, its compliance with the Metadata Guideline, and a
list of suggested actions. Before describing the interface, we answer one question: \emph{who should
use it}, the curator (an EarthPortal administrator), the author of an ontology, or both?

% =====================================================================
\section{Development}

\subsection{Who does the work?}

Each suggested action of our audits says who can take it: the \textbf{author} (in the ontology file or
the submission form), the \textbf{curator} (the EarthPortal team), or the \textbf{portal} itself (values
it generates). For the 68 artefacts of the snapshot of 16 October 2026, the actions of priorities 1 to 3
(Mandatory properties, and properties that cost O'FAIRe points) are split as follows:

\begin{center}\small\renewcommand{\arraystretch}{1.2}
\begin{tabular}{@{}l>{\raggedright\arraybackslash}p{8.4cm}rr@{}}
\toprule
\textbf{Who} & \textbf{Examples of properties} & \textbf{Actions} & \textbf{Share} \\
\midrule
author  & License, Creator, Modification date, Other identifier (DOI), Version IRI & 2,626 & 96\,\% \\
curator & Categories, Groups, Curator, Curation date, Included in catalog & 81 & 3\,\% \\
portal  & CSV dump, Download URL, URI lookup endpoint & 37 & 1\,\% \\
\midrule
        & \textbf{Total} & \textbf{2,744} & \\
\bottomrule
\end{tabular}
\captionof{table}{Who can take the actions of priorities 1 to 3, for the 68 artefacts (16 October 2026).}
\label{tab:who}
\end{center}

Most of the information that is missing is known only by the authors: the licence they chose, who
created the ontology, its DOI. A curator cannot invent it. So an interface used by curators alone would
mostly produce messages to authors; and an interface open to anyone would let people change metadata
that is not theirs.

\begin{keypoint}
Our answer is \textbf{both, with different rights, and in two steps}. First, a curator interface,
reserved to EarthPortal administrators. Then, once the suggestions have been validated by the team, an
author view limited to the ontologies each author administers.
\end{keypoint}

OntoPortal already knows these two roles: an \textbf{administrator of the portal} can edit every
ontology, and each ontology has its own administrators (the field \prop{administeredBy}), who can edit
its metadata. The interface should reuse these rights, not create new ones.

\noindent\begin{minipage}{\linewidth}
\begin{center}\small\renewcommand{\arraystretch}{1.25}
\begin{tabular}{@{}>{\raggedright\arraybackslash}p{2.6cm}>{\raggedright\arraybackslash}p{6.0cm}>{\raggedright\arraybackslash}p{6.0cm}@{}}
\toprule
& \textbf{Curator (portal administrator)} & \textbf{Author (administrator of an ontology)} \\
\midrule
Sees & the whole catalogue & only the ontologies they administer \\
Acts on & the curator actions; the checks of the portal; the requests sent to authors & the author actions of their own ontologies \\
Decides & accept, edit, reject (with a reason) or postpone each suggestion & fill or correct their metadata, or explain why not \\
When & step 1 & step 2, after the team has validated the suggestions \\
\bottomrule
\end{tabular}
\captionof{table}{The two roles of the interface.}
\label{tab:roles}
\end{center}
\end{minipage}


\subsection{Four principles}

\begin{enumerate}
  \item \textbf{The pipeline proposes, a person decides.} Nothing is written to the portal without a
    click. Our checks only see whether a value is \emph{present}, not whether it is \emph{correct}: for
    example, 55 artefacts have the same default value for their object description properties.
  \item \textbf{Every decision is recorded}: who, when, which suggestion, and why. A rejected
    suggestion, with its reason, is not proposed again at the next run. A change can be undone.
  \item \textbf{Both views of quality are shown, but the score does not decide alone.} Each action
    says why the guideline asks for it and what it is worth in O'FAIRe. A property such as \emph{Abstract}
    costs no O'FAIRe point, but the guideline recommends it.
  \item \textbf{Every result has a date.} The audits describe one snapshot of the catalogue. The
    interface shows that date, and marks a suggestion as outdated when the value has changed since.
\end{enumerate}

\subsection{Three screens}

\begin{enumerate}
  \item \textbf{Catalogue}: one row per ontology, with its O'FAIRe and reachable scores, its guideline
    compliance (for example ``Mandatory 16/20, Recommended 2/57''), its open actions, and the date of
    its last audit. It is our \texttt{index.md}, made interactive (sorting, filters).
  \item \textbf{Ontology}: the two parts of its audit, then its actions grouped by who can take them.
    For each action, the current value (read live from the portal), the suggestion, the evidence, and
    the buttons to decide.
  \item \textbf{Property}: one property of the guideline across the whole catalogue, for example the
    23 artefacts without a Creator. This view shows patterns; a group action is still confirmed
    ontology by ontology.
\end{enumerate}

Figure~\ref{fig:screen} sketches the second screen, the one where most decisions are taken.

\begin{figure}[H]
\centering
\begin{tikzpicture}[font=\scriptsize,
  box/.style={draw=dtgrey,rounded corners=1.5pt,inner sep=4pt},
  btn/.style={draw=dtnavy,rounded corners=1.5pt,inner sep=2pt,font=\tiny,text=dtnavy}]
  \draw[dtgrey,rounded corners=3pt] (0,0) rectangle (15,8.4);
  \fill[dtnavy] (0,7.8) rectangle (15,8.4);
  \node[anchor=west,text=white,font=\small\bfseries] at (0.3,8.1) {STY: Semantic Types Ontology};
  \node[anchor=east,text=white] at (14.7,8.1) {snapshot 2026-10-16 \quad artefact};
  \node[box,anchor=north west,text width=6.6cm] at (0.3,7.5)
    {\textbf{1. FAIRness (O'FAIRe)}\\ 39/100 \quad reachable 49.5\,\% \quad rank 68 of 68\\
     largest losses: I1Q1 (18), R1.1Q1 (15), F1Q2 (11)};
  \node[box,anchor=north west,text width=6.6cm] at (7.8,7.5)
    {\textbf{2. Metadata Guideline}\\ Mandatory 16/20 \quad Recommended 2/57\\ Optional 0/20 \quad
     \textcolor{dtnavy}{show the 99 properties}};
  \node[anchor=west,font=\small\bfseries,text=dtnavy] at (0.3,5.75) {3. Actions};
  \node[anchor=west] at (2.3,5.75) {\textbf{author} \quad curator \quad portal};
  \foreach \x/\h in {0.3/Priority, 1.3/Property, 4.4/Guideline, 6.7/O'FAIRe, 10.1/Value now}
    \node[anchor=west,text=dtgrey,font=\tiny] at (\x,5.35) {\h};
  \foreach \y/\p/\name/\why/\pts/\val in {
     5.0/P1/License (4.2)/Mandatory/15 pts (R1.1Q1)/{\textit{empty}},
     4.35/P1/Creator (5.2)/Mandatory/2 pts (R1.2Q1)/{\textit{empty}},
     3.7/P2/Version IRI (1.9)/Recommended/9 pts (F1Q4)/{\textit{empty}},
     3.05/P3/Other identifier (1.14)/Optional/{17 pts (A1Q1, F1Q2)}/{\textit{empty}},
     2.4/P4/Abstract (2.1)/Recommended/no O'FAIRe point/{\textit{empty}}} {
    \draw[dtgrey!50] (0.3,\y-0.32) -- (14.7,\y-0.32);
    \node[anchor=west,font=\scriptsize\bfseries] at (0.3,\y) {\p};
    \node[anchor=west] at (1.3,\y) {\name};
    \node[anchor=west] at (4.4,\y) {\why};
    \node[anchor=west] at (6.7,\y) {\pts};
    \node[anchor=west] at (10.1,\y) {\val};
    \node[btn] at (11.7,\y) {accept};
    \node[btn] at (12.6,\y) {edit};
    \node[btn] at (13.45,\y) {reject};
    \node[btn] at (14.3,\y) {later};
  }
  \node[anchor=west,text=dtgrey] at (0.3,1.65) {\ldots};
  \node[box,anchor=south west,text width=14.0cm,fill=dtband] at (0.3,0.25)
    {\textbf{Request to the author}: the author actions above become one polite message, ready to send
     to the contacts of the ontology, with a link to its audit page.};
\end{tikzpicture}
\caption{A sketch of the ``Ontology'' screen, with the data of STY on 16 October 2026. The buttons
\emph{accept} and \emph{edit} write to the portal only for the curator actions; for the author actions,
the curator prepares a request instead.}
\label{fig:screen}
\end{figure}

\subsection{How the interface gets its data}

The interface does not need to run our Python code. Since v0.3.0, the pipeline writes \textbf{one JSON
file per ontology} (\prop{audits/json/<ACRONYM>.json}); the audit page that people read is written from
the same file, so the two always agree. Each action looks like this:

\begin{lstlisting}
{ "priority": 3, "why": "Optional, costing O'FAIRe points", "who_fixes": "author",
  "id": "1.14", "entry": "Other identifier", "usage": "Optional",
  "portal_attribute": "dct:identifier", "questions": ["A1Q1", "F1Q2"], "points_share": 17,
  "note": "For F1Q2, only a registered DOI gives the 11 points; any other valid URL gives 6." }
\end{lstlisting}

The JSON holds no value of the metadata, and so no name or e-mail address. The interface reads the
current values live from the portal API, with the rights of the person who is logged in, and writes
the accepted changes through the same API.

\begin{figure}[H]
\centering
\begin{tikzpicture}[font=\small, node distance=0.9cm,
  b/.style={draw=dtnavy,rounded corners=2pt,align=center,minimum height=1.1cm,text width=3.3cm,fill=dtband},
  a/.style={->,dtnavy,line width=0.8pt}]
  \node[b] (p) {Pipeline\\ \scriptsize notebooks 01 to 05};
  \node[b,right=of p] (j) {One JSON file\\ per ontology};
  \node[b,right=of j] (u) {Curation\\ interface};
  \node[b,below=0.8cm of u] (api) {EarthPortal API\\ \scriptsize current values and changes};
  \node[b,below=0.8cm of j] (md) {Audit pages\\ \scriptsize for people};
  \draw[a] (p) -- (j); \draw[a] (j) -- (u); \draw[a] (j) -- (md);
  \draw[a,<->] (u) -- (api);
  \draw[a,dashed] (api.south) -- ++(0,-0.45) -| node[pos=0.25,below,font=\scriptsize] {next snapshot} (p.south);
\end{tikzpicture}
\caption{The data flow. The pipeline and the interface only share the JSON files.}
\label{fig:flow}
\end{figure}

\subsection{Preventing is better than curing}

The submission form of EarthPortal (Figure~\ref{fig:form}) asks for name, acronym, visibility,
administrators, URI, description, representation language, status, location, date of creation and
contacts. It does not ask for the \textbf{licence}, the \textbf{creator} or the \textbf{modification
date}, and it offers the \textbf{categories} without requiring them. These are exactly four of the five
Mandatory properties most often empty in the catalogue:

\begin{center}\small\renewcommand{\arraystretch}{1.2}
\begin{tabular}{@{}lrl@{}}
\toprule
\textbf{Mandatory property} & \textbf{Artefacts where it is empty} & \textbf{In the submission form} \\
\midrule
Modification date & 27 & not asked \\
Location          & 24 & empty by design when a file is uploaded \\
Creator           & 23 & not asked \\
Categories        & 20 & offered, not required \\
License           & 15 & not asked \\
\bottomrule
\end{tabular}
\captionof{table}{The Mandatory properties most often empty among the 68 artefacts (16 October 2026).}
\label{tab:form}
\end{center}

Asking for these properties at submission time would prevent many actions from ever reaching the
curation interface. A small detail on the same form: the text at the top of step 2 ends with an
untranslated \texttt{\%\{link\}}.

\begin{figure}[H]
\centering
\includegraphics[height=7.2cm]{EarthPortal_Ontology-Submission-production_Step1.png}\hfill
\includegraphics[height=7.2cm]{EarthPortal_Ontology-Submission-production_Step2.png}\hfill
\includegraphics[height=3.6cm]{EarthPortal_Ontology-Submission-production_Step3.png}
\caption{The three steps of the submission form of EarthPortal (production, 29 October 2026).}
\label{fig:form}
\end{figure}

% =====================================================================
\section{Conclusion}

We suggest one curation interface with two roles: first for the EarthPortal curators, then for the
authors, each on their own ontologies, using the rights OntoPortal already has. It reads one JSON file
per ontology produced by our pipeline, never writes without a person's decision, and keeps a record of
every decision. In parallel, adding four properties to the submission form would reduce the work at the
source. I would appreciate your feedback, in particular on the two roles and on the place of the
interface: inside the EarthPortal administration pages, or as a separate tool.

\end{document}
